\chapter{Zależności zewnętrzne i poziom samowystarczalności}
\label{app:wyniki-zewnetrzne}

W skrypcie dowód oznacza rzeczywiste
wyprowadzenie ze wskazanych wcześniejszych
faktów. Jeśli korzystamy z długiego
twierdzenia spoza tekstu, podajemy
hipotezy, miejsce użycia i źródło.
Poniższy katalog ułatwia rozpoznanie
tych punktów zależności.

\section{Wyniki wejściowe i głębokie lematy}
\label{sec:katalog-wynikow}

\begin{description}[leftmargin=2.9cm,style=nextline]
\item[Niezmienniczość obszaru.]
  Ciągła injekcja z otwartego
  $U\subset\R^n$ do $\R^n$
  ma otwarty obraz. Użyta w
  Rozdziale~\ref{ch:rozmaitosci}
  do niezależności wymiaru i
  brzegu od map.
  Źródło:
  \href{https://pi.math.cornell.edu/~hatcher/AT/AT.pdf}
  {A.\ Hatcher, \emph{Algebraic Topology},
  twierdzenie 2B.3}.
  Alternatywny dowód samej
  niezmienniczości wymiaru rozmaitości
  przez homologię lokalną jest poniżej.

\item[Twierdzenie o funkcji odwrotnej.]
  Dla $C^1$ odwzorowania między
  otwartymi podzbiorami $\R^n$
  z odwracalną pochodną w punkcie
  istnieje lokalna odwrotność $C^1$;
  dla danych $C^\infty$ jest ona
  $C^\infty$. Użyte przy mapach,
  zanurzeniach i funkcjach Morse'a.
  Źródło z dowodem:
  \href{https://www.ocw.mit.edu/courses/18-101-analysis-ii-fall-2005/resources/lecture6/}
  {V.\ Guillemin, \emph{Analysis II},
  wykład 6}.

\item[Triangulacja gładka.]
  Każda gładka rozmaitość
  drugoprzeliczalna dopuszcza
  triangulację zgodną ze
  strukturą gładką; w przypadku
  zwartym może być skończona.
  Używamy jej do przejścia od
  geometrii do skończonych
  kompleksów łańcuchowych.
  Źródło:
  \href{https://www.ams.org/bull/1935-41-08/S0002-9904-1935-06140-3/S0002-9904-1935-06140-3.pdf}
  {S.~S.\ Cairns, \emph{Triangulation of the
  Manifold of Class One}, §1--2}.
  Pełny dowód nie jest krótkim
  uzupełnieniem tego skryptu.

\item[Kohomologia Čech i singularna.]
  Każde pokrycie rozmaitości ma
  rozdrobnienie będące dobrym pokryciem.
  Dobre pokrycia tworzą więc rodzinę
  kofinalną w definicji kohomologii
  Čecha. Dla dobrego pokrycia kompleks
  Čecha oblicza kohomologię nerwu,
  a lemat o nerwie utożsamia jego
  typ homotopijny z rozmaitością.
  Dla współczynników w dyskretnej
  grupie abelowej otrzymujemy stąd
  izomorfizm kohomologii Čecha
  i singularnej. Interpretacja
  kocyklu przejść struktury
  spinowej w Rozdziale~\ref{ch:grass-spin}
  używa tego porównania;
  tam, gdzie wystarcza stopień
  pierwszy, podano też argument
  przez monodromię pętli.
  Źródło szczegółowego porównania:
  \href{https://pi.math.cornell.edu/~hatcher/AT/AT.pdf}
  {A.\ Hatcher, \emph{Algebraic Topology},
  wniosek 4G.3 o nerwie dobrego pokrycia
  oraz §2.1 i §3.1 o kohomologii
  symplicjalnej i singularnej}.

\item[Trik Whitneya i ruchy uchwytów.]
  Wysokowymiarowe usuwanie par
  punktów przecięcia wymaga
  ściągalności okręgu Whitneya
  w dopełnieniu, warunków na wymiar
  i obramowanego dysku Whitneya.
  Użyte w
  Rozdziałach~\ref{ch:hcobordyzm}--\ref{ch:przeszkoda-chirurgiczna}.
  Źródło warunków i konstrukcji:
  \href{https://www.math.columbia.edu/~jmorgan/Lecture_IIIA_hcobordism_Contd.pdf}
  {J.\ Morgan, \emph{The $h$-Cobordism Theorem},
  wykład IIIA}.

\item[Twierdzenia Walla.]
  Przy normalnym odwzorowaniu
  stopnia jeden w wymiarze
  co najmniej $6$, z prostą
  równoważnością na brzegu,
  zanik przeszkody w $L_n^s$
  jest równoważny możliwości
  chirurgii do prostej
  równoważności. Rozdział~\ref{ch:przeszkoda-chirurgiczna}
  rozpisuje rachunek przecięć i konstrukcję
  bazy hiperbolicznej nad $\Z[\pi]$,
  sprawdza zgodność obramowania w pojedynczej
  chirurgii poniżej środka, przeprowadza
  skończoną redukcję poniżej środka w zamkniętym
  przypadku parzystym oraz wyprowadza dwa
  ciągi dokładne jąder.
  Rozpisano też ciąg dokładny, dualność
  i znoszenie przecięć wyznaczające prosty
  lagranżjan z lematu~5.7 Walla. Kryterium
  osadzania środkowej sfery z warunku
  $\mu=0$ rozwinięto do parowania samoprzecięć,
  konstrukcji pętli i dysku Whitneya oraz
  kontroli ram normalnych. Produktowe
  otoczenie dysku i lokalny ruch obramowany
  nadal korzystają z
  \href{https://webhomes.maths.ed.ac.uk/~v1ranick/surgery/hcobord.pdf}
  {twierdzenia~6.6 i lematów~6.7, 6.13 Milnora}.
  Osobnymi wejściami pozostają przygotowanie kobordyzmu,
  bazowość i prosta dualność jego jąder,
  realizacja względnych immersji oraz
  pełna realizacja ruchów formacji
  nieparzystych. W szczególności rachunek
  lagranżjanu nie zastępuje geometrycznych
  hipotez twierdzeń~1.3 i~1.4 ani
  lematów~2.4, 2.6 i~5.5 Walla.
  Źródło:
  \href{https://webhomes.maths.ed.ac.uk/~v1ranick/books/scm.pdf}
  {C.~T.~C. Wall, \emph{Surgery on Compact
  Manifolds}, twierdzenia 1.2, 5.2,
  5.6, 6.4 i lemat 5.7; dla ciągu
  chirurgii także §10}.
  Rozdziały~\ref{ch:przeszkoda-chirurgiczna}
  oraz~\ref{ch:ciag-chirurgii}
  rozpisują ich zastosowanie,
  lecz nie zastępują wszystkich
  lematów Walla.

\item[Triangulacja powierzchni i wymiar pięć.]
  Dowód klasyfikacji zwartych powierzchni
  w Twierdzeniu~\ref{thm:classification-surfaces-24}
  korzysta z istnienia skończonej triangulacji
  topologicznej powierzchni; źródło:
  \href{https://pi.math.cornell.edu/~hatcher/Papers/TorusTrick.pdf}
  {A.\ Hatcher, \emph{The Kirby Torus Trick for Surfaces}}.
  Osobny przypadek gładkich
  $5$-rozmaitości wynika z
  \href{https://www.sas.rochester.edu/mth/sites/doug-ravenel/otherpapers/barden.pdf}
  {D.\ Barden, \emph{Simply Connected
  Five-Manifolds}, twierdzenie 2.2};
  jego założenia zapisano w
  Twierdzeniu~\ref{thm:Barden-5-24}.

\item[Sygnatura i sfery Milnora.]
  Dla zamkniętej zorientowanej
  $8$-rozmaitości $C$ używamy
  $\operatorname{sign}(C)=
  \langle(7p_2-p_1^2)/45,[C]\rangle$.
  Źródło:
  \href{https://webhomes.maths.ed.ac.uk/~v1ranick/papers/hirzrem.pdf}
  {F.\ Hirzebruch, \emph{The Signature Theorem},
  wzór na $L_2$}.
  Wzory charakterystyczne
  wiązek nad $S^4$ pochodzą
  z \href{https://sites.math.rutgers.edu/~feehan/teaching/math866/milnor7sphere.pdf}
  {J.\ Milnor, \emph{On Manifolds Homeomorphic
  to the 7-Sphere}, lemat 3}.

\item[Analityczna teoria Hodge'a.]
  Dla laplasjanu Hodge'a na formach zwartej
  rozmaitości riemannowskiej bez brzegu
  jądro jest skończenie wymiarowe, a każda
  gładka forma prostopadła do jądra ma
  jednoznaczny gładki przeciwobraz
  prostopadły do jądra. Jest to
  Twierdzenie~\ref{thm:hodge-analytic-input},
  użyte w Rozdziale~\ref{ch:hodge-bochner}
  do istnienia rozkładu; jego konsekwencje
  algebraiczne i wzór Bochnera tam dowodzimy.
  Źródło:
  \href{https://doi.org/10.2307/1969552}
  {K.\ Kodaira, \emph{Harmonic Fields in
  Riemannian Manifolds} (1949)}.

\item[Analityczna teoria potoku Ricciego.]
  Gładki ściśle paraboliczny układ
  quasiliniowy drugiego rzędu na zwartej
  rozmaitości bez brzegu ma dla gładkich
  danych lokalnie w czasie jednoznaczne
  gładkie rozwiązanie; używamy także
  odpowiedniej wersji dla równania ciepła
  odwzorowań z metryką dziedziny zależną
  od czasu. W
  Twierdzeniu~\ref{thm:ricci-short-time}
  te fakty stosujemy po poprawce DeTurcka.
  Wyprowadzenie głównego symbolu i cofnięcie
  metryki przez dyfeomorfizmy są w rozdziale;
  same oszacowania paraboliczne są zewnętrzne.
  Źródło twierdzenia o krótkim czasie:
  \href{https://projecteuclid.org/journals/journal-of-differential-geometry/volume-17/issue-2/Three-manifolds-with-positive-Ricci-curvature/10.4310/jdg/1214436922.full}
  {R.\ S.\ Hamilton, \emph{Three-manifolds with
  positive Ricci curvature} (1982), \S\,4}.
  Źródło konstrukcji i równania cechowania:
  \href{https://projecteuclid.org/journals/journal-of-differential-geometry/volume-18/issue-1/Deforming-metrics-in-the-direction-of-their-Ricci-tensors/10.4310/jdg/1214509286.full}
  {D.\ DeTurck, \emph{Deforming Metrics in the
  Direction of Their Ricci Tensors} (1983)}.

\item[Topologiczna dualność i wygładzanie PL w wymiarze cztery.]
  W Rozdziale~\ref{ch:freedman-formy} formę przecięcia
  zamkniętej rozmaitości TOP opisujemy przez topologiczną
  dualność Poincarégo. Używamy także wyniku, że
  struktura PL w wymiarze cztery dopuszcza zgodne
  wygładzenie. Rachunki macierzy uchwytów i formy $E_8$
  są w rozdziale; te dwa fakty kategorialne pozostają
  wejściami zewnętrznymi. Źródło:
  \href{https://archive.mpim-bonn.mpg.de/4789/2/FreedmanQuinn-TopologyOf4Manifolds-Reformatted2013.pdf}
  {M.\ H.\ Freedman i F.\ Quinn, \emph{Topology of
  4-Manifolds}, \S\,8.7 i \S\,10.2}.

\item[Osadzanie dysków i uchwyty Cassona.]
  W Rozdziale~\ref{ch:uchwyty-cassona} dowodzimy
  własności skończonych pięter, ale przejście do
  parami rozłącznych lokalnie płaskich dysków
  przy hipotezach o grupie podstawowej i sferach
  dualnych jest wejściem zewnętrznym:
  Twierdzenie~\ref{thm:casson-disk-embedding},
  \href{https://archive.mpim-bonn.mpg.de/4789/2/FreedmanQuinn-TopologyOf4Manifolds-Reformatted2013.pdf}
  {Freedman--Quinn, \emph{Topology of 4-Manifolds},
  Twierdzenie 5.2}. Konstrukcję geometrycznie
  dualnych sfer w konkluzji uzupełnia
  \href{https://arxiv.org/abs/2006.05209}
  {M.\ Powell, A.\ Ray i P.\ Teichner,
  \emph{The 4-dimensional disc embedding theorem
  and dual spheres}, Twierdzenie A}; w pierwotnym
  dowodzie Freedmana--Quinna brakowało tego kroku.
  Standardowość topologiczna
  uchwytu Cassona jako pary to osobne wejście:
  Twierdzenie~\ref{thm:casson-homeomorphism},
  \href{https://projecteuclid.org/journals/journal-of-differential-geometry/volume-17/issue-3/The-topology-of-four-dimensional-manifolds/10.4310/jdg/1214437136.full}
  {Freedman, \emph{The topology of four-dimensional
  manifolds} (1982), Twierdzenie 1.1}.

\item[Klasyfikacja Freedmana i niezmiennik Kirby'ego--Siebenmanna.]
  Twierdzenie~\ref{thm:freedman-classification}
  o zamkniętych zorientowanych jednospójnych
  rozmaitościach TOP, warunek
  $\operatorname{ks}\equiv\sigma/8\pmod2$ dla
  form parzystych oraz
  wynik o przeszkodzie do stabilnego wygładzenia
  są wejściami zewnętrznymi w
  Rozdziale~\ref{ch:freedman-klasyfikacja}.
  Źródła:
  \href{https://archive.mpim-bonn.mpg.de/4789/2/FreedmanQuinn-TopologyOf4Manifolds-Reformatted2013.pdf}
  {Freedman--Quinn, \emph{Topology of 4-Manifolds},
  \S\,10.1--10.2} oraz
  \href{https://projecteuclid.org/journals/journal-of-differential-geometry/volume-17/issue-3/The-topology-of-four-dimensional-manifolds/10.4310/jdg/1214437136.full}
  {Freedman, \emph{The topology of four-dimensional
  manifolds} (1982), Twierdzenie 1.5}.
\end{description}

\section{Wybrany dowód: wymiar z homologii lokalnej}
\label{sec:local-dimension-appendix}

W Rozdziale~\ref{ch:rozmaitosci}
wyprowadziliśmy jednoznaczność
wymiaru z twierdzenia Brouwera.
Po wprowadzeniu homologii
można otrzymać inny dowód,
bez zakładania samej
niezmienniczości obszaru.

\begin{stwierdzenie}[Homologiczna niezmienniczość
wymiaru rozmaitości]
\label{prop:dimension-local-homology}
Jeżeli niepuste rozmaitości
topologiczne bez brzegu
$M^n,N^m$ są homeomorficzne,
to $n=m$.
\end{stwierdzenie}
\begin{proof}
Dla $p\in M$ wybieramy
mapę z małą kulą
$U\cong\R^n$. Przez
wycinanie i długi ciąg
pary mamy
\[
 H_i(M,M\setminus\{p\};\Z)
 \cong H_i(\R^n,\R^n\setminus\{0\};\Z)
 \cong\widetilde H_{i-1}(S^{n-1};\Z).
\]
Drugi izomorfizm wynika z
ściągalności $\R^n$
i radialnego cofnięcia
$\R^n\setminus\{0\}$
na $S^{n-1}$.
Dla $n\geq1$ grupa
po lewej jest $\Z$
dokładnie w stopniu $i=n$
i zerowa w pozostałych
stopniach. Dla $n=0$
lokalna grupa $H_0$
jest $\Z$, co daje
ten sam sposób odczytu.
Homeomorfizm par
$(M,M\setminus\{p\})$
i $(N,N\setminus\{h(p)\})$
zachowuje wszystkie
grupy lokalne, zatem
stopień ich jedynej
niezerowej grupy jest
taki sam: $n=m$.
\end{proof}
